% answers_short.tex —— 2016 年·全国Ⅱ卷·理科 · 简短答案速查
% 由 tools/make_short_answers.py 自 answers.tex 的 \gksolhead/\gkanswer 提取，
% 并与 2016/2016_数学真题_全国2理_答案要点.md 逐题比对后生成，勿手改（改源 MD 后重跑）。

\begin{gknotice}
\textbf{使用说明}：本节只给最终答案便于快速核对，完整推导见解析版。答案为编写者独立推导，\textbf{非官方标准答案}。
\end{gknotice}

\gkpart{选择题}
\begin{center}\begin{tabular}{ccccccccccccc}
\toprule
题号 & 1 & 2 & 3 & 4 & 5 & 6 & 7 & 8 & 9 & 10 & 11 & 12 \\
\midrule
答案 & A & C & D & A & B & C & B & C & D & C & A & B \\
\bottomrule
\end{tabular}\end{center}

\gkpart{填空题}
\begin{itemize}[leftmargin=2.6em]
  \item[\textbf{13.}] \(\frac{21}{13}\)
  \item[\textbf{14.}] ②③④
  \item[\textbf{15.}] 甲拿 \(\{1,3\}\)
  \item[\textbf{16.}] \(b=1-\ln2\)
\end{itemize}

\gkpart{解答题}
\begin{itemize}[leftmargin=2.6em]
  \item[\textbf{17.}] (1) \(b_1=0,\ b_{11}=1,\ b_{101}=2\)；(2) \(1893\)
  \item[\textbf{18.}] (1) \(0.55\)；(2) U；(3) \(\frac{123}{100}=1.23\)
  \item[\textbf{19.}] (1) \(D'H\perp\) 平面 \(ABCD\)；(2) \(\sin\) 二面角 \(B-D'-A-C=\frac{2\sqrt{95}}{25}\)
  \item[\textbf{20.}] (1) \(\frac{144}{49}\)；(2) \(\sqrt[3]{2}<k<2\)
  \item[\textbf{21.}] (1) \(f\) 在 \((-\infty,-2)\)、\((-2,+\infty)\) 上分别严格递增；当 \(x>0\) 时，\((x-2)e^x+x+2>0\)；(2) \(\left(\frac12,\frac{e^2}{4}\right]\)
\end{itemize}

\gkpart{选做题}
\begin{itemize}[leftmargin=2.6em]
  \item[\textbf{22.}] (1) \(B,C,G,F\) 四点共圆；(2) \(\frac12\)
  \item[\textbf{23.}] (1) \(\rho^2+12\rho\cos\theta+11=0\)；(2) 斜率为 \(\pm\frac{\sqrt{15}}3\)
  \item[\textbf{24.}] (1) \(M=(-1,1)\)；(2) \(|a+b|<|1+ab|\)
\end{itemize}
